\documentclass[10pt,landscape]{article}

\usepackage{alltt}
\usepackage{array}
\usepackage{booktabs}
\usepackage{enumitem}
\usepackage[T1]{fontenc}
\usepackage[landscape,margin=13mm,footskip=1pt,includefoot]{geometry}
\usepackage{graphicx}
\usepackage{hyperref}
\usepackage{multicol}
\usepackage[small,compact]{titlesec}
\usepackage{upquote}
\usepackage{xcolor}
\usepackage{xspace}

\graphicspath{{figs/}}

\pagestyle{empty}
\parindent=0pt

\setitemize{itemsep=0.1mm,parsep=1pt,leftmargin=10pt}
\setenumerate{itemsep=0.2mm,parsep=1pt}

\hypersetup{
    colorlinks=true,        % false: boxed links; true: colored links
    linkcolor=red,          % color of internal links
    citecolor=cyan,         % color of links to bibliography
    filecolor=magenta,      % color of file links
    urlcolor=blue           % color of external links
}

\def\todo#1{\textit{\textcolor{red}{TODO: #1}}}
\newcommand{\minisec}[1]{\textbf{#1}\\}

\begin{document}

\begin{multicols*}{3}
  
{\Huge\scshape Bro Cheat Sheet}

\hfill\includegraphics[width=.4\linewidth]{bro-logo-small}
\vspace{-44pt}

{\scriptsize
\url{http://www.bro-ids.org}\\
\texttt{\href{mailto:info@bro-ids.org}{info@bro-ids.org}}\\
\url{https://github.com/broids/cheat-sheet}\\
Version: \today
}
\vspace{5pt}

\section*{Startup}

\texttt{bro \textit{[options] [file \dots]}}\\
\texttt{\textit{file}} \dotfill Bro policy script or \texttt{stdin}\\
\texttt{-e \textit{code}} \dotfill Augment policies by given code\\
\texttt{-h} \dotfill Display command line options\\
\texttt{-i \textit{iface}} \dotfill Read from given interface\\
\texttt{-p \textit{pfx}} \dotfill Add given prefix to policy resolution\\
\texttt{-r \textit{file}} \dotfill Read from given PCAP file\\
\texttt{-w \textit{file}} \dotfill Write to given file in PCAP format\\
\texttt{-x \textit{file}} \dotfill Print contents of state file\\
\texttt{-C} \dotfill Ignore invalid checksum\\
\texttt{-Z} \dotfill Generate documentation for all loaded scripts\\

\vspace{-10pt}

\section*{Language}

In this document, lowercase letters represent instance variables and uppercase
letters represent types. In general, \texttt{x} is an instance of type
\texttt{T} and \texttt{y} and instance of type \texttt{U}. Argument names and
record fields begin begin with \texttt{a}, \texttt{b}, \ldots, and \texttt{z}
represents a default instance variable which takes on the type of the
right-hand side expression. Parenthesis after a type, e.g., \texttt{T()},
denote constructor invocation.

\subsection*{Expressions}

%\minisec{Operators}
\verb|!| \dotfill Negation\\
\verb|$|, \verb|?$| \dotfill Dereference, record field existence\\
\verb|+|, \verb|-|, \verb|*|, \verb|/|, \verb|%| \dotfill Arithmetic\\
\verb|++|, \verb|--| \dotfill Post-increment, post-decrement\\
\verb|+=|, \verb|-=|, \verb|*=|, \verb|/=| \dotfill Arithmetic and assignment\\
\verb|==|, \verb|!=| \dotfill Equality, inequality\\
\verb|<|, \verb|<=|, \verb|>=|, \verb|>|
  \dotfill Less/greater than (or equal)\\
\verb|&&|, \verb#||# \dotfill Conjunction, disjunction\\
\verb|in|, \verb|!in| \dotfill Membership or pattern matching\\
\verb|[x]| \dotfill Index strings and containers\\
\verb#|x|# \dotfill Cardinality/size for strings and containers\\
\verb|::| \dotfill Access \texttt{module} or \texttt{enum} namespace\\
\texttt{\textit{expr}$\;$?$\;$\textit{expr}$\;$:$\;$\textit{expr}}
  \dotfill Ternary if-then-else

\subsection*{Statements}

Basic statement \dotfill \texttt{\textit{stmt};}\\ 
Code block \dotfill \texttt{\{ \textit{stmt}; ... \}}\\ 
Script import \dotfill \texttt{@load \textit{path}}\\
Type declaration \dotfill \texttt{type~\textit{name}:~T}\\
Print expression to \texttt{stdout} \dotfill
\texttt{print \textit{expr}}\\

\minisec{Variables}
Scope qualifier \dotfill \texttt{local}, \texttt{global}\\ 
Declaration \dotfill \texttt{\textit{scope} x:~T}\\
Declaration \& Definition \dotfill \texttt{\textit{scope} x = T()}\\
Assignment \dotfill \texttt{x = T()}\\

\minisec{Constants}
Constant qualifier \dotfill \texttt{const} \\
Redefinition \dotfill \texttt{redef x \textit{op} T()}\\

\minisec{Namespaces}
\texttt{module \textit{ns}} \dotfill Set current namespace to \texttt{ns}\\
\verb|export { ... }| \dotfill Export global symbols of this file \\

\minisec{Functions}
Declaration \dotfill \verb|function f(a: T, ...): R|\\
Invocation \dotfill \verb|f(x, ...)|\\
Anonymous \dotfill \texttt{z = function(...):~R~\{ ... \}}\\

\minisec{Events}
Signature \dotfill \verb|event e(a: T, ...)|\\
Queuing \dotfill \verb|event e(x, ...)|\\
Scheduling \dotfill \verb|schedule 10 secs { e(x, ...) }|\\

\vspace{-10pt}
%\minisec{Control Flow}
\begin{multicols*}{3}
\textbf{Branching}
\begin{alltt}
if (\textit{expr})
    \dots
else if (\textit{expr})
    \dots
else
    \dots
\end{alltt}

\textbf{Iteration}
\begin{alltt}
for (i in x)
    \dots
\end{alltt}

\textbf{Conditional}
\begin{alltt}
when (x = y)
    \dots
\end{alltt}
\vspace{-10pt}

\begin{alltt}
\textrm{\textbf{Control}}\vspace{5pt}
break
continue
next
return
\end{alltt}
\end{multicols*}

\subsection*{Types}

\minisec{Basic}
\verb|addr| \dotfill IP address (\verb|127.0.0.1|)\\
\verb|bool| \dotfill Boolean flag (\verb|T|, \verb|F|)\\
\verb|count| \dotfill 64-bit unsigned integer (\verb|42|)\\
\verb|double| \dotfill Double-precision floating point (\verb|99.9|)\\
\verb|int| \dotfill 64-bit signed integer (\verb|-7|)\\
\verb|interval| \dotfill Time interval
  (\verb|8 sec|/\verb|min|/\verb|hr|/\verb|day[s]|)\\
\verb|pattern| \dotfill Regular expression (\verb|/^br[oO])$/|)\\
\verb|port| \dotfill Transport-layer port
  (\verb|22/tcp|, \verb|53/udp|)\\
\verb|string| \dotfill String of bytes (\verb|"foo"|)\\
\verb|subnet| \dotfill CIDR subnet mask (\verb|10.0.0.0/8|)\\
\verb|time| \dotfill Absolute epoch time (\verb|1320977325|)\\

\minisec{Enumerables}
Declaration \dotfill \verb|enum { FOO, BAR }|\\
Construction \dotfill \texttt{\textit{scope}} \verb|x = FOO|\\

\minisec{Records}
Declaration \dotfill \verb|record { a: T, b: U, ... }|\\
Construction \dotfill \texttt{\textit{scope}} \verb|r = [$a=x, $b=y, ...]|\\
Access \dotfill \verb|z = r$a|\\
Assignment \dotfill \verb|r$b = y|\\
Deletion \dotfill \verb|delete r$a|\\

\minisec{Sets}
Declaration \dotfill \verb|set[T]|\\
Constructor \dotfill \verb|set()|\\
Construction \dotfill \texttt{\textit{scope}} \verb|s = { x, ... }|\\
Access \dotfill \verb|z = s[x]|\\
Addition \dotfill \verb|add s[x]|\\
Deletion \dotfill \verb|delete s[x]|\\

\minisec{Tables}
Declaration \dotfill \verb|table[T] of U|\\
Constructor \dotfill \verb|table()|\\
Construction \dotfill \texttt{\textit{scope}} \verb|t = { [x] = y, ... }|\\
Access \dotfill \verb|z = t[x]|\\
Assignment \dotfill \verb|add t[x] = y|\\
Deletion \dotfill \verb|delete t[x]|\\

\minisec{Vectors}
Declaration \dotfill \verb|vector of T|\\
Constructor \dotfill \verb|vector()|\\
Construction \dotfill \texttt{\textit{scope}} \verb|v = { x, ... }|\\
Access \dotfill \verb|z = v[0]|\\
Assignment \dotfill \verb|v[42] = x|\\
% FIXME: not yet supported in Bro
%Deletion \dotfill \verb|delete v[3]|
\end{multicols*}

\begin{multicols*}{2}
\subsection*{Attributes}
Attributes occur at the end of type/event declarations and change their
behavior. The syntax is \verb|&key| or \verb|&key=val|, e.g.,
\verb|type T: set[count] &read_expire=5min| or
\verb|event foo() &priority=-3|.\\

\verb|&optional| \dotfill Allow record field to be missing\\
\verb|&default=x| \dotfill Use default value when element is not in container\\
\verb|&redef| \dotfill Make constants redefinable\\
% Users should use the logging framework instead.
%\verb|&rotate_interval=x| \dotfill Rotate file after time \verb|x|\\
%\verb|&rotate_size=x| \dotfill Rotate file after reaching size limit \verb|x|\\
\verb|&add_func=f| \dotfill Call \verb|f| \todo{after?} adding element to
  container\\
\verb|&delete_func=f| \dotfill Call \verb|f| right before deleting element from
  container \\
\verb|&expire_func=f| \dotfill Call \verb|f| right before container element 
  expires\\
\verb|&read_expire=x| \dotfill Remove element after not reading it for time
  \verb|x|\\
\verb|&write_expire=x| \dotfill Remove element after not writing it for time
  \verb|x|\\
\verb|&create_expire=x| \dotfill Remove element after time \verb|x| from
  insertion\\
\verb|&persistent| \dotfill Write state to disk (per default on shutdown)\\
\verb|&synchronized| \dotfill Synchronize variable across nodes\\
% TODO: these are probably less relevant for users.
%\verb|&postprocessor| \dotfill \\
%\verb|&encrypt| \dotfill \\
%\verb|&match| \dotfill \\
\verb|&raw_output| \dotfill Do not escape non-ASCII characters when writing to
a file\\
\verb|&mergeable| \dotfill Prefer set union to assignment for synchronized
  state\\
\verb|&priority=x| \dotfill Execution priority of event handler\\
\verb|&group="x"| \dotfill Events in the same group can be
activated/deactivated together \\
\verb|&log| \dotfill Write record field to log \\

\subsection*{Built-In Functions (BIFs)}
\linespread{0.9}

\subsubsection*{Core}

\begin{itemize}
  \item \verb|length(v: any): count|\\
    Returns the number of elements in the container \texttt{v}.
  \item \verb|same_object(o1: any, o2: any): bool|\\
    Check whether \texttt{o1} and \texttt{o2} reference the same internal
    object.
  \item \verb|clear_table(v: any): any|\\
    Remove all elements from the set or table \texttt{v}.
\end{itemize}

\subsubsection*{String Processing}

\begin{itemize}
  \item \verb|byte_len(s: string): count|\\
    Returns the number of characters (i.e., bytes) of \texttt{s}.
  \item \verb|sub_bytes(s: string, start: count, n: int): string|\\
    Get a substring of \texttt{s}, starting at position \texttt{start} and
    having length \texttt{n}.
  \item \verb|split(s string, re: pattern): table[count] of string|\\
    Split \texttt{s} into an array using \texttt{re} to separate the elements.
  \item \verb|split1(s string, re: pattern): table[count] of string|\\
    Same as \texttt{split}, but \texttt{s} is only split once (if possible) at
    the earliest position and an array of two strings is returned. An array of
    one string is returned when \texttt{s} cannot be split.
  \item \verb|split_all(s: string, re: pattern): table[count] of string|\\
    Same as \texttt{split}, but also include the matching separators, e.g.,
    \verb|split_all("a-b--cd", /(\-)+/)| returns
    \verb|{"a", "-", "b", "--", "cd"}|. Odd-indexed elements do not match the
    pattern and even-indexed ones do.
  \item
\begin{verbatim}
split_n(s: string, re: pattern, incl_sep: bool,
        max_num_sep: count): table[count] of string
\end{verbatim}
    Similar to \verb|split1| and \verb|split_all|, but \verb|incl_sep|
    indicates whether to include matching separators and \verb|max_num_sep| the
    number of times to split \texttt{s}.
%  \item 
%\begin{verbatim}
%function split_complete(s string, re: pattern,
%                        other: set[string], incl_sep: bool,
%                        max_num_sep: count): table[count] of string
%\end{verbatim}
%    Same as \verb|split_n|, except that the extra parameter \texttt{other} is
%    not used.
  \item \verb|sub(s: string, re: pattern, repl: string): string|\\
    Substitutes \texttt{repl} for the first occurrence of \texttt{re} in
    \texttt{s}.
  \item \verb|gsub(s string, re: pattern, repl: string): string|\\
    Same as \texttt{sub} except that all occurrences of \texttt{re} are
    replaced.
  \item \verb|strcmp(s1: string, s2: string): int|
    Lexicographically compare \texttt{s1} and \texttt{s2}. Returns an integer
    greater than, equal to, or less than 0 according as \texttt{s1} is greater
    than, equal to, or less than \texttt{s2}.
  \item \verb|strstr(big: string, little: string): count|\\
    Locate the first occurence of \texttt{litle} in \texttt{big}.
    Returns 0 if \texttt{little} is not found in \texttt{big}.
  \item \verb|subst_string(s: string, from: string, to: string): string|\\
    Substitute each (non-overlapping) appearance of \texttt{from} in
    \texttt{s} to \texttt{to}, and return the resulting string.
  \item \verb|to_lower(s: string): string|\\
    Returns a copy of \texttt{s} with each letter converted to lower case.
  \item \verb|to_upper(s: string): string|\\
    Returns a copy of \texttt{s} with each letter converted to upper case.
  \item \verb|clean(s: string): string|\\
    Replace non-printable characters in \texttt{s} with escaped sequences,
    with the mappings
    \verb|0| $\rightarrow$ \verb|\0|, 
    \verb|DEL| $\rightarrow$ \verb|^?|,
    values $\le 26$ $\rightarrow$ \verb|^[A-Z]|,
    and values not in $[32, 126]$~$\rightarrow$~\verb|%XX|.
  \item \verb|to_string_literal(s: string): string|\\
    Same as clean, but with different mappings: 
    values not in $[32, 126]$~$\rightarrow$~\verb|%XX|,
    \verb|\|~$\rightarrow$~\verb|\\|, 
    \verb|'|~$\rightarrow$~\verb|\'|,
    \verb|"|~$\rightarrow$~\verb|\"|.
  \item \verb|is_ascii(s: string): bool|\\
    Returns false if any byte value of \texttt{s} is greater than 127, and true
    otherwise.
  \item \verb|escape_string(s: string): string|\\
    Returns a printable version of \texttt{s}. Same as \texttt{clean} except
    that non-printable characters are removed.
  \item \verb|string_to_ascii_hex(s: string): string|\\
    Returns an ASCII hexadecimal representation of a string.
  \item \verb|str_split(s: string, idx: vector of count): vector of string|\\
    Splits \texttt{s} into substrings, taking all the indices in
    \texttt{idx} as cutting points; \texttt{idx} does not need to be sorted,
    and can have multiple entries. Out-of-bounds indices are ignored.
  \item \verb|strip(s: string): string|\\
    Strips whitespace at both ends of \texttt{s}.
  \item \verb|string_fill(len: int, source: string): string|\\
    Generates a string of size \texttt{len} and fills it with repetitions of
    \texttt{source}.
  \item \verb|str_shell_escape(source: string): string|\\
    Takes a string and escapes characters that would allow execution of
    commands at the shell level. Must be used before including strings in
    \texttt{system()} or similar calls.
  \item \verb|find_all(s: string, re: pattern) : string_set|\\
    Returns all occurrences of \texttt{re} in \texttt{s} (or an empty empty set
    if none).
  \item \verb|find_last(s: string, re: pattern) : string|\\
    Returns the last occurrence of \texttt{re} in \texttt{s}. If not found,
    returns an empty string.  Note that this function returns the match that
    starts at the largest index in the string, which is not necessarily the
    longest match.  For example, a pattern of \texttt{/.*/} will return the
    final character in the string.
  \item \verb|hexdump(data: string) : string|\\
    Returns a hex dump for \texttt{data}. The hex dump renders 16 bytes per
    line, with hex on the left and ASCII (where printable) on the right. Based
    on Netdude's hex editor code.
\end{itemize}

\end{multicols*}

\end{document}
